\documentclass[10pt,a4paper]{article}

\usepackage[a4paper,margin=1.8cm,columnsep=0.7cm]{geometry}
\usepackage[T1]{fontenc}
\usepackage[utf8]{inputenc}
\usepackage{lmodern}
\usepackage{microtype}
\usepackage{graphicx}
\usepackage{booktabs}
\usepackage{array}
\usepackage{enumitem}
\usepackage{caption}
\usepackage{multicol}
\usepackage{amsmath}
\usepackage{amssymb}
\usepackage{hyperref}

\hypersetup{
  colorlinks=true,
  linkcolor=black,
  citecolor=black,
  urlcolor=blue,
  pdfauthor={Chaoui Kouraichi Ahmed, Huang Ruisi},
  pdftitle={Making LLMs Faster: Reusable Memory for Documents and Images}
}

\setlength{\parindent}{0pt}
\setlength{\parskip}{0.45em}
\setlength{\columnsep}{0.7cm}
\setlist[itemize]{leftmargin=*,topsep=0.2em,itemsep=0.1em}
\setlist[enumerate]{leftmargin=*,topsep=0.2em,itemsep=0.1em}
\captionsetup{font=small,labelfont=bf}

\newcommand{\reporttitle}{MAKING LLMS FASTER}
\newcommand{\reportsubtitle}{Reusable Memory for Documents and Images}
\newcommand{\reportdate}{June 2026}

\newcommand{\authorblock}[3]{%
  \begin{minipage}[t]{0.42\textwidth}
    \centering
    #1\\
    #2\\
    #3
  \end{minipage}
}

\begin{document}

\begin{center}
  {\Large\bfseries \reporttitle\par}
  \vspace{0.25em}
  {\large \reportsubtitle\par}
  \vspace{1.4em}

  \authorblock{Chaoui Kouraichi Ahmed}{EURECOM, Sophia Antipolis, France}{\texttt{chaoui@eurecom.fr}}
  \hfill
  \authorblock{Huang Ruisi}{EURECOM, Sophia Antipolis, France}{\texttt{Ruisi.Huang@eurecom.fr}}

  \vspace{1em}
  \reportdate
\end{center}

\begin{abstract}
  Large language models are increasingly used as semantic operators in
  data-processing systems that operate over documents, images, and other
  unstructured inputs. However, repeatedly applying LLMs to the same base items
  across filters, mappings, joins, and alternative execution plans introduces a
  high computational cost. This project studies reusable Key-Value (KV) cache
  compression as a way to reduce redundant inference work while preserving task
  quality. We compare Expected Attention Press, KVzip, Finch with context
  prefill tuning (CPT), and Finch without CPT across text and vision workloads,
  with a focus on classification and extraction tasks. The goal is to understand
  how different methods select tokens to retain, how they trade off quality,
  latency, and memory footprint, and when the additional cost of query-dependent
  compression is justified by better downstream performance.
\end{abstract}

\begin{multicols}{2}

\section{Introduction}

Large language models (LLMs) have become a practical building block for data
systems that need to process unstructured and multimodal data. Unlike
traditional database operators, which usually rely on exact matching or
predefined schemas, LLM-based operators can answer semantic questions over raw
documents and images. Recent LLM-native systems, such as Palimpzest and Lotus,
make this idea explicit: users describe high-level operations such as filtering
documents by meaning, extracting attributes, or classifying items, and the
system translates these operations into LLM calls.

This shift creates new opportunities, but also a central systems problem. In a
semantic pipeline, the same base item may be reused many times: a document can
be filtered, mapped, joined, or evaluated under several alternative execution
plans. If each operator sends the full input back through the model, the system
recomputes a large amount of identical context. This redundancy becomes
especially costly for long documents and visual inputs, where the number of
tokens can be large and model inference is both memory-intensive and slow.

Transformer-based LLMs already use Key-Value (KV) caching during inference to
avoid recomputing attention states for previous tokens. For each layer, the
model stores keys and values that summarize the processed context and reuses
them when generating new tokens. KV caching is essential for efficient
generation, but it also introduces a memory cost that grows with the context
length and the number of layers. In workloads where many base items must be
kept available for repeated queries, this memory footprint becomes a limiting
factor.

KV-cache compression addresses this limitation by keeping only a subset of the
cached tokens. Query-agnostic methods, such as Expected Attention Press and
KVzip, select tokens that appear globally important during the prefill phase.
Their main advantage is reuse: once a compact cache has been produced for a
document or image, it can be stored and shared across many future queries.
Query-dependent methods, such as Finch, instead select tokens according to the
incoming query. In this report, we evaluate Finch both with context prefill
tuning (CPT) and without CPT in order to isolate the effect of query-guided
selection from the additional cache preparation step.

The objective of this project is to evaluate this tradeoff in the context of
LLM-native data systems. We study how existing KV-cache compression strategies
decide which tokens to keep, compare query-agnostic and query-dependent methods
including Expected Attention Press, KVzip, Finch with CPT, and Finch without
CPT, and analyze the practical conditions under which each family is
preferable. In particular, we are interested in whether query-dependent
compression can improve quality without increasing the memory footprint, for
example by storing several aggressively compressed caches instead of one
less-compressed reusable cache. The report focuses on classification and
extraction tasks over text and vision inputs, and uses these experiments to
reason about the broader latency, memory, and reuse tradeoffs of compressed LLM
memory.

\section{Related Work}

\subsection{LLM-Native Data Systems}

Recent work on LLM-native data systems studies how large language models can be
integrated into data-processing pipelines as reusable computational operators.
Systems such as Palimpzest~\cite{palimpzest} and Lotus~\cite{lotus} expose
high-level semantic operators for tasks such as filtering, mapping,
classification, extraction, and joining over unstructured data. Instead of
requiring users to manually write one prompt for each item and each operation,
these systems organize LLM calls as part of a larger query plan.

This view is useful for document and image workloads because many operations
depend on semantic understanding rather than exact string matching. For
example, a system may need to determine whether a document expresses a
sentiment, whether an image contains an object relevant to a query, or whether a
piece of text contains a specific attribute. In such settings, the LLM behaves
like a semantic operator inside a dataflow. However, the same base item is often
processed repeatedly across different operators or optimizer choices. This
repetition motivates mechanisms that can reuse intermediate model computation
rather than recomputing the same document or image representation for every
query.

\subsection{KV-Cache Reuse and Compression}

Transformer models rely on Key-Value (KV) caches to avoid recomputing attention
states for previously processed tokens during generation. For long inputs, the
cache can represent a large fraction of inference memory. This issue becomes
more pronounced in LLM-native data systems, where many documents or images may
need to be cached and reused. KV-cache compression therefore aims to keep the
most useful tokens while reducing the memory footprint of the stored cache.

Query-agnostic methods compress an input without knowing the future query.
Expected Attention Press estimates which tokens are likely to receive attention
from future queries and retains those tokens in the cache~\cite{expected-attention}.
KVzip follows the same reusable-cache motivation, but uses context
reconstruction to identify compact caches that preserve the useful information
of the original context~\cite{kvzip}. These approaches are attractive for
systems workloads because the compressed cache can be computed once and reused
across many downstream tasks. Their limitation is that they must preserve
information that may be useful for many possible queries, which can make them
less specialized for any single query.

\subsection{Query-Dependent Compression}

Query-dependent compression takes a different approach: it uses the incoming
query to decide which cached tokens are relevant. Finch is a representative
method in this family, using prompt guidance to select cache entries that are
important for the specific request~\cite{finch}. This can improve quality under
tight memory budgets because the compressed cache is optimized for the query at
hand. The cost is reduced reusability: if the query changes, the selected cache
may no longer be appropriate and additional computation may be needed.

This report compares both families in the same experimental setting. Expected
Attention Press and KVzip represent query-agnostic reusable compression, while
Finch with CPT and Finch without CPT represent query-dependent compression. The
comparison is designed to clarify not only which method gives better task
quality, but also when the extra query-specific computation is justified by
improvements in accuracy, F1 score, latency, or memory efficiency.

\section{Experimental Setup}

\subsection{Tasks and Datasets}

We evaluate KV-cache compression across three experiment lines covering
both text and vision modalities.

\textbf{Movie Reviews (Text, Sentiment Classification).}
We use a corpus of 1{,}000 English-language movie reviews for binary
sentiment classification (positive/negative). Each review is processed
with Qwen2.5-0.5B-Instruct. The task tests whether a compressed KV cache
retains enough semantic information to determine the overall sentiment of
a document. Compression ratios tested are
$r \in \{0.2, 0.4, 0.6, 0.8\}$.

\textbf{Movie Reviews (Text, Information Extraction).}
A subset of 120 reviews per query is used for three extraction tasks:
sentiment polarity extraction (query~010), movie title extraction
(query~011), and praised-actor extraction (query~012). These tasks test
whether the compressed cache preserves fine-grained factual details that
the model must retrieve verbatim.

\textbf{Email Dataset (Text, Filter + Extract).}
We evaluate on an email corpus of 100 emails per query, covering 20
queries: 10 filter queries (boolean yes/no, e.g.\ ``Does this email
discuss a meeting?'') and 10 extraction queries (e.g.\ ``Extract the
sender's department''). The model is Qwen2.5-0.5B-Instruct. Compression
ratios tested are $r \in \{0.2, 0.4, 0.6, 0.8\}$.

\textbf{Artwork (Vision-Language).}
We use 10 paintings from a curated art dataset, each paired with 20
queries (10~filter, 10~extract). The model is LLaVA-Next 8B
(\texttt{llava-hf/llama3-llava-next-8b-hf}). Compression ratios tested
are $r \in \{0.4, 0.8, 0.95\}$. This experiment line tests whether
KV-cache compression is viable for multimodal inputs where visual patch
tokens dominate the cache.

\subsection{Evaluation Metrics}

We report the following metrics across all experiments:

\begin{itemize}
  \item \textbf{Accuracy}: fraction of correct predictions (used for
    sentiment classification and extraction tasks).
  \item \textbf{Macro F1}: harmonic mean of precision and recall, averaged
    across classes. This is the primary metric for classification tasks
    and accounts for class imbalance.
  \item \textbf{Per-query Precision, Recall, F1}: for filter tasks, we
    compute binary precision, recall, and F1 per query against
    ground-truth positive sets. For extraction tasks, we compute macro
    P/R/F1 across answer categories.
  \item \textbf{Mean F1}: for the artwork benchmark, we average F1 across
    all 20 queries to produce a single summary score per method and
    compression ratio.
  \item \textbf{Mean latency (ms)}: wall-clock time per inference,
    measured end-to-end including cache loading and generation.
\end{itemize}

All filter queries in the email dataset achieve perfect precision (1.0)
across both methods and all compression ratios, indicating that compressed
caches never produce false positives. The quality variation is driven
entirely by recall.

\begin{center}
  \centering
  \fbox{\rule{0pt}{3.2cm}\rule{0.9\columnwidth}{0pt}}
  \captionof{figure}{Experiment overview: datasets, models, methods, and
  compression ratios used across the three experiment lines.}
  \label{fig:dataset-binary}
\end{center}

\begin{center}
  \centering
  \fbox{\rule{0pt}{4cm}\rule{0.9\columnwidth}{0pt}}
  \captionof{figure}{Distribution of filter vs.\ extraction queries across
  the email and artwork benchmarks, showing the number of records
  evaluated per query type.}
  \label{fig:wide-distribution}
\end{center}

\section{Methods}

\subsection{Expected Attention Press}

Expected Attention (EA)~\cite{expected-attention} is a query-agnostic,
score-based compression method. During the prefill phase, EA assigns an
importance score to each KV pair by estimating the expected attention it
will receive from future queries.

The method models the distribution of future queries using the statistics
(mean and covariance) of the queries observed during prefill. It then
computes, in closed form, how much attention each key is expected to
receive from a typical future query. Tokens with the highest expected
attention are retained; the rest are discarded according to the
compression ratio. Sink tokens (the first 4 positions) are always
protected from pruning.

EA requires only a single forward pass and adds negligible overhead to
prefill, making it the fastest method in our comparison.

\subsection{KVzip}

KVzip~\cite{kvzip} is a query-agnostic method that uses context
reconstruction to score token importance. Unlike EA, which estimates
importance statistically, KVzip asks the model to reconstruct the
original context in chunks and measures which tokens receive the most
cross-attention during this reconstruction process.

This approach is more computationally expensive---it requires multiple
additional forward passes (typically 3--5 reconstruction chunks)---but
it measures importance more directly by observing how the model actually
uses each token. The resulting cache is still query-agnostic and
reusable. In our experiments, KVzip is approximately $4\times$ slower
than EA in terms of cache preparation latency.

\subsection{Finch without CPT}

Finch~\cite{finch} is a query-dependent compression method that uses the
incoming prompt to guide token selection. It extends the SnapKV approach
with two modifications: (1)~a dynamic observation window whose size
adapts based on delimiter token positions in the input, and (2)~key
re-rotation after pruning to restore continuous positional embeddings
(RoPE).

In the variant without context prefill tuning (CPT), Finch compresses the
KV cache purely based on the attention pattern between the query window
and the context tokens. The observation window size is set to 16 tokens.

To ensure a fair comparison with query-agnostic methods, we adjust the
compression ratio to account for the window tokens that Finch always
retains, so that after removing window tokens the same fraction of context
tokens is kept as in the other methods.

\subsection{Finch with CPT}

Finch with context prefill tuning (CPT) extends the base Finch method by
using a pre-computed, document-specific cached prompt to guide
compression. For each base document, a short textual summary or
representative prompt is stored alongside the KV cache. At compression
time, this prompt is used instead of (or in addition to) the incoming
query to determine which tokens are important.

The motivation is to combine the quality benefits of query-dependent
selection with partial reusability: if the cached prompt captures the
essential content of the document, the resulting compressed cache may
transfer across similar queries. In practice, however, the CPT
preparation adds latency and storage overhead (one prompt per document),
and the quality improvements depend on how well the cached prompt
anticipates future queries.

\subsection{Compression and Cache-Reuse Settings}

We test compression ratios $r \in \{0.2, 0.4, 0.6, 0.8\}$ for text
workloads and $r \in \{0.4, 0.8, 0.95\}$ for vision workloads, where
$r$ denotes the fraction of KV pairs removed. A ratio of $r=0.8$ means
that 80\% of the original tokens are discarded and only 20\% are
retained.

For query-agnostic methods (EA and KVzip), the compressed cache is
computed once per document during an offline prefill phase and stored to
disk. At inference time, the pre-compressed cache is loaded and reused
for each incoming query without any recomputation. This models the
intended use case in LLM-native data systems, where a document may be
queried hundreds of times.

For query-dependent methods (Finch with and without CPT), the cache must
be recomputed or adapted for each new query. We measure this additional
cost in the latency numbers reported in our results.

All experiments use the \texttt{kvpress}
library\footnote{\url{https://github.com/NVIDIA/kvpress}} for
compression and \texttt{KeyRerotationPress} wrapping where applicable to
maintain continuous positional embeddings after token pruning.

\section{Results}

\subsection{Overall Performance}

Table~\ref{tab:main-results} summarizes the key results across all
experiment lines. The most striking finding is the robustness of
query-agnostic methods: EA and KVzip maintain near-baseline quality even
at aggressive compression ratios ($r=0.8$), while Finch variants degrade
significantly on the vision workload.

\begin{center}
  \centering
  \captionof{table}{Summary of main results across workloads. Reviews
    report macro F1 (n=1000, no CPT). Artwork reports mean F1 across 20
    queries. Extract reports macro F1 at $r=0.8$.}
  \label{tab:main-results}
  \begin{tabular}{lcccc}
    \toprule
    Method & Reviews & Artwork & Artwork & Extract \\
           & $r{=}0.8$ & $r{=}0.4$ & $r{=}0.95$ & $r{=}0.8$ \\
    \midrule
    EA              & 0.822 & 0.621 & 0.631 & 0.517 \\
    KVzip           & 0.811 & 0.639 & 0.646 & 0.568 \\
    Finch (no CPT)  & 0.781 & 0.534 & 0.393 & --- \\
    Finch + CPT     & 0.770 & 0.407 & 0.259 & --- \\
    \bottomrule
  \end{tabular}
\end{center}

\subsection{Text Workloads}

\textbf{Sentiment classification (reviews).}
Table~\ref{tab:reviews-results} shows that EA and KVzip achieve nearly
identical accuracy and macro F1 across all compression ratios, with
negligible degradation from $r=0.2$ to $r=0.8$. EA is consistently
$\sim$4$\times$ faster than KVzip ($\sim$395\,ms vs.\
$\sim$1{,}730\,ms per inference). Enabling CPT does not improve quality
for either method and reduces F1 by 5--9 points.

\begin{center}
  \centering
  \captionof{table}{Sentiment classification on 1{,}000 reviews
    (Qwen2.5-0.5B, no CPT).}
  \label{tab:reviews-results}
  \begin{tabular}{lcccc}
    \toprule
    Ratio & EA Acc & KVzip Acc & EA F1 & KVzip F1 \\
    \midrule
    0.2 & 0.847 & 0.843 & 0.820 & 0.816 \\
    0.4 & 0.852 & 0.838 & 0.826 & 0.810 \\
    0.6 & 0.846 & 0.850 & 0.819 & 0.824 \\
    0.8 & 0.848 & 0.839 & 0.822 & 0.811 \\
    \bottomrule
  \end{tabular}
\end{center}

\begin{center}
  \centering
  \fbox{\rule{0pt}{3.2cm}\rule{0.9\columnwidth}{0pt}}
  \captionof{figure}{EA vs KVzip accuracy and macro F1 on 1{,}000 movie
    reviews across compression ratios.}
  \label{fig:reviews-metrics}
\end{center}

\textbf{Email filter tasks.}
On the email dataset, both methods achieve perfect precision (1.0) across
all 10 filter queries and all compression ratios. Recall varies by query
difficulty (from 0.37 to 0.98), but the mean F1 across filter queries
remains stable at $\sim$0.84 regardless of compression ratio. This
confirms that filter tasks---which only need to detect the
\emph{presence} of relevant information---are highly robust to KV-cache
compression.

\textbf{Email extraction tasks.}
Extraction tasks show dramatically lower accuracy. Two queries
(query\_012 and query\_017) achieve 0.0 accuracy for \emph{both} EA and
KVzip at \emph{all} compression ratios, indicating that the required
factual tokens are systematically pruned by global scoring methods.
KVzip slightly outperforms EA on extraction (mean accuracy $\sim$0.27
vs.\ $\sim$0.21 at $r=0.8$), but both methods struggle. EA accuracy
degrades from 0.26 at $r=0.2$ to 0.21 at $r=0.8$, while KVzip remains
more stable at $\sim$0.27.

\textbf{Movie review extraction.}
Table~\ref{tab:extract-results} shows per-query extraction results.
Query~012 (praised actor) is an outlier where EA strongly outperforms
KVzip (0.867 vs.\ 0.658), while KVzip is better on sentiment and title
extraction.

\begin{center}
  \centering
  \captionof{table}{Movie review extraction tasks
    (Qwen2.5-0.5B, $r{=}0.8$).}
  \label{tab:extract-results}
  \begin{tabular}{lcc}
    \toprule
    Query & EA F1 & KVzip F1 \\
    \midrule
    010 (sentiment)     & 0.517 & 0.568 \\
    011 (movie title)   & 0.417 & 0.454 \\
    012 (praised actor) & 0.867 & 0.658 \\
    \bottomrule
  \end{tabular}
\end{center}

\begin{center}
  \centering
  \fbox{\rule{0pt}{4cm}\rule{0.9\columnwidth}{0pt}}
  \captionof{figure}{EA vs KVzip macro F1 on extraction tasks across
    compression ratios.}
  \label{fig:extract-f1}
\end{center}

\subsection{Vision Workloads}

KVzip achieves the highest mean F1 on the artwork benchmark (0.639 at
$r=0.4$, stable at 0.646 for $r \geq 0.8$), followed closely by EA
(0.621--0.632). Both query-agnostic methods are remarkably stable across
compression ratios.

\begin{center}
  \centering
  \captionof{table}{Artwork benchmark mean F1 (LLaVA-Next 8B, 20
    queries, 10 paintings).}
  \label{tab:artwork-results}
  \begin{tabular}{lccc}
    \toprule
    Method & $r{=}0.4$ & $r{=}0.8$ & $r{=}0.95$ \\
    \midrule
    EA              & 0.621 & 0.633 & 0.631 \\
    KVzip           & 0.639 & 0.646 & 0.646 \\
    Finch (no CPT)  & 0.534 & 0.344 & 0.393 \\
    Finch + CPT     & 0.407 & 0.329 & 0.259 \\
    \bottomrule
  \end{tabular}
\end{center}

Finch without CPT starts at a moderate 0.534 at $r=0.4$ but degrades
sharply to 0.344 at $r=0.8$. Adding CPT makes the situation worse:
Finch with CPT drops from 0.407 to 0.259 across the same range. This
suggests that query-dependent compression is particularly fragile for
multimodal inputs, where visual patch tokens form a dense, redundant
representation that does not benefit from query-specific filtering.

\begin{center}
  \centering
  \fbox{\rule{0pt}{3.2cm}\rule{0.9\columnwidth}{0pt}}
  \captionof{figure}{Mean F1 across 20 artwork queries for four
    compression methods at ratios 0.4, 0.8, and 0.95.}
  \label{fig:artwork-f1}
\end{center}

\section{Discussion}

\subsection{Quality, Latency, and Memory Tradeoffs}

Our results reveal a central empirical finding: \textbf{quality is
remarkably stable across compression ratios for query-agnostic methods.}
On sentiment classification, EA achieves F1 of 0.820 at $r=0.2$ and
0.822 at $r=0.8$---effectively no degradation despite removing
$4\times$ more tokens. On artwork, KVzip achieves 0.639 at $r=0.4$ and
0.646 at $r=0.95$. This suggests that the vast majority of tokens in
the KV cache contribute very little to the model's downstream
predictions, and that the scoring functions used by EA and KVzip
successfully identify the small set of tokens that matter.

\textbf{Latency.}
EA is consistently $\sim$4$\times$ faster than KVzip across all
workloads ($\sim$395\,ms vs.\ $\sim$1{,}730\,ms on reviews,
$\sim$3{,}000\,ms vs.\ $\sim$4{,}100\,ms on emails). This is because
EA requires only a single forward pass, while KVzip runs multiple
reconstruction passes. Finch with query-dependent compression adds
further latency overhead at low compression ratios
($\sim$6{,}800\,ms at $r=0.2$) due to the dynamic window computation.

\begin{center}
  \centering
  \fbox{\rule{0pt}{3.2cm}\rule{0.9\columnwidth}{0pt}}
  \captionof{figure}{Mean latency per inference for EA vs KVzip on
    1{,}000 reviews across compression ratios.}
  \label{fig:latency}
\end{center}

\textbf{Memory.}
At $r=0.8$, the compressed KV cache is $5\times$ smaller than the
uncompressed cache. For a system managing thousands of documents, this
translates directly into the number of documents that can be held in
GPU memory simultaneously, which is the key bottleneck for LLM-native
data systems like Palimpzest.

\subsection{When to Use Query-Agnostic or Query-Dependent Compression}

Our experiments identify a clear practical boundary between the two
compression families.

\textbf{Query-agnostic methods (EA, KVzip) are preferable when:}
\begin{itemize}
  \item The cache will be reused across multiple queries. Since the
    compressed cache is computed once during offline prefill, the
    amortized cost decreases with each additional query.
  \item The task is classification or filtering. These tasks require
    detecting the \emph{presence} of semantic features, which correlates
    well with globally important tokens.
  \item Latency is a concern. EA in particular adds near-zero overhead
    to prefill and achieves quality comparable to KVzip at
    $\sim$4$\times$ lower latency.
\end{itemize}

\textbf{Query-dependent methods (Finch) are preferable when:}
\begin{itemize}
  \item The task requires extracting specific factual details (names,
    numbers, locations) that may not rank highly in global importance
    scores.
  \item The cache is used for a single query and will not be reused,
    making the recomputation cost acceptable.
  \item The compression ratio is moderate ($r \leq 0.4$). At higher
    compression, Finch degrades significantly, especially on vision
    tasks.
\end{itemize}

\textbf{Limitations observed.}
Finch with CPT consistently underperforms Finch without CPT on both text
and vision workloads. This suggests that the cached prompt does not
adequately capture the information needs of future queries, and the
additional preparation overhead is not justified. On the artwork
benchmark, Finch + CPT drops to a mean F1 of 0.259 at $r=0.95$,
compared to 0.646 for KVzip at the same ratio---a gap of nearly 0.4 F1
points.

Two extraction queries (query\_012 and query\_017 on the email dataset)
achieve 0.0 accuracy for \emph{all} methods at \emph{all} compression
ratios. These represent a hard failure mode where the required factual
tokens are not among the globally highest-scoring KV pairs and are
systematically pruned regardless of the method used.

\section{Conclusion and Future Work}

This project evaluated four KV-cache compression methods---Expected
Attention Press, KVzip, Finch without CPT, and Finch with CPT---across
text and vision workloads in the context of LLM-native data systems.

Our main findings are:
\begin{enumerate}
  \item \textbf{Query-agnostic methods are remarkably robust.} EA and
    KVzip maintain near-baseline quality even when 80\% of KV pairs are
    discarded, across both text sentiment classification (F1 $\sim$0.82)
    and vision filter tasks (mean F1 $\sim$0.64). Quality is stable
    across compression ratios.
  \item \textbf{Filter tasks are robust; extraction tasks are fragile.}
    All methods achieve perfect precision and stable recall on boolean
    filter queries. Extraction tasks, by contrast, suffer dramatic
    accuracy drops---including complete failure on certain
    queries---because the required factual tokens are pruned by global
    scoring.
  \item \textbf{EA Pareto-dominates KVzip on the quality--latency
    frontier.} EA achieves F1 scores within 1--2 points of KVzip while
    running $\sim$4$\times$ faster. For systems workloads where latency
    matters, EA is the preferred choice.
  \item \textbf{Query-dependent compression (Finch) does not help on
    vision.} Finch degrades sharply on the artwork benchmark, and adding
    CPT worsens performance further. Query-dependent selection appears
    ill-suited for the dense token structure of visual inputs.
\end{enumerate}

\textbf{Future work.} Several directions emerge from these results:
\begin{itemize}
  \item \textbf{Task-adaptive compression}: automatically detecting the
    task type (filter vs.\ extract) and adjusting the compression ratio
    accordingly---aggressive for filters, conservative for extraction.
  \item \textbf{Multi-cache strategies}: storing multiple aggressively
    compressed query-dependent caches per document (e.g., one for
    sentiment, one for entity extraction) instead of one moderately
    compressed query-agnostic cache, to achieve higher coverage within
    the same memory budget.
  \item \textbf{Hybrid methods}: using EA for most transformer layers
    but applying KVzip-style reconstruction scoring to a small number
    of critical layers to balance quality and latency.
  \item \textbf{Larger models}: scaling to 7B--70B models, which may
    exhibit sharper attention distributions and tolerate more aggressive
    compression.
  \item \textbf{Theoretical analysis}: formalizing the relationship
    between attention score distributions and compression robustness to
    provide principled guidance for method selection and budget
    allocation.
\end{itemize}

\begin{thebibliography}{99}

\bibitem{finch}
Giulio Corallo and Paolo Papotti.
\newblock Finch: Prompt-guided key-value cache compression for large language
models.
\newblock \emph{Transactions of the Association for Computational Linguistics},
12:1517--1532, 2024.

\bibitem{expected-attention}
Alessio Devoto, Maximilian Jeblick, and Simon Jegou.
\newblock Expected attention: KV cache compression by estimating attention from
future queries distribution.
\newblock arXiv preprint, 2025.

\bibitem{kvzip}
Jang-Hyun Kim, Jinuk Kim, Sangwoo Kwon, Jae W. Lee, Sangdoo Yun, and Hyun Oh
Song.
\newblock KVzip: Query-agnostic KV cache compression with context
reconstruction.
\newblock 2025.

\bibitem{palimpzest}
Chunwei Liu, Matthew Russo, Michael Cafarella, Lei Cao, Peter Baile Chen, Zui
Chen, Michael Franklin, Tim Kraska, Samuel Madden, Rana Shahout, and Gerardo
Vitagliano.
\newblock Palimpzest: Optimizing AI-powered analytics with declarative query
processing.
\newblock In \emph{Proceedings of the Conference on Innovative Data Systems
Research (CIDR)}, 2025.

\bibitem{lotus}
Liana Patel, Siddharth Jha, Melissa Pan, Harshit Gupta, Parth Asawa, Carlos
Guestrin, and Matei Zaharia.
\newblock Semantic operators: A declarative model for rich, AI-based data
processing.
\newblock 2025.

\end{thebibliography}

\end{multicols}

\section*{Appendix: Theoretical Framework for KV Cache Compression}

This appendix establishes a rigorous mathematical foundation for the empirical successes and failure modes of KV cache compression algorithms.

\subsection*{A.1 The Attention Mass Concentration Theorem}

\textbf{Motivation:} Empirical benchmarks show that dropping up to 80\% of KV cache tokens (compression ratio $r=0.8$) results in negligible degradation in F1 scores for filter tasks.

\textbf{Definition A.1 (Attention Mass):} Let $s_1, s_2, \dots, s_n$ be the attention importance scores of $n$ context tokens. Assuming a compression algorithm retains the top $k$ tokens (where $k = n(1-r)$), the retained attention mass is:
\begin{equation}
M(r) = \frac{\sum_{i \in \text{top-}k} s_i}{\sum_{i=1}^{n} s_i}
\end{equation}

\textbf{Lemma A.1 (Pre-softmax Logit Distribution):} In a well-trained transformer, the dot product attention logits $z_i = \frac{q \cdot k_i}{\sqrt{d}}$ for a given query $q$ across a diverse key distribution tend toward an Extreme Value Distribution (e.g., Gumbel) due to the maximization objective inherent in cross-entropy training. Let the PDF of the logits be $f(z) \sim e^{-(z - \mu)} e^{-e^{-(z - \mu)}}$.

\textbf{Theorem A.1 (Attention Mass Concentration):} 
If the pre-softmax logits follow a Gumbel-like extreme value distribution, then the post-softmax attention weights $s_i = \frac{e^{z_i}}{\sum_j e^{z_j}}$ follow a power-law (Zipfian) distribution $P(s > x) \sim x^{-\alpha}$ in the tail.
Under this power-law distribution with shape parameter $\alpha > 1$, the attention mass lost by discarding the bottom fraction $r$ of tokens is bounded by:
\begin{equation}
1 - M(r) = \mathcal{O}\left(r^{\frac{\alpha - 1}{\alpha}}\right)
\end{equation}

\textit{Proof Sketch:}
The transformation $s_i \propto e^{z_i}$ from a Gumbel tail $P(Z > z) \sim e^{-z}$ yields a Pareto distribution. Integrating the lowest $r$ fraction of the probability density gives the mass of the discarded tokens, yielding exactly $r^{1 - 1/\alpha}$. For empirical LLM layers, $\alpha \approx 2$ to $4$, so the discarded mass vanishes rapidly.

\subsection*{A.2 The Task Hardness Separation Theorem}

\textbf{Motivation:} While filter tasks survive $r=0.8$ compression seamlessly, specific extraction tasks suffer a complete accuracy collapse (dropping to 0.0) at the same compression ratio.

\textbf{Theorem A.2 (Phase Transition of Pruning Fragility):}
Let $S_k$ be the set of $k$ tokens retained by the compression scoring function. Assume the correlation between the scoring function and the true relevance is $\rho \in (0, 1]$.
\begin{enumerate}
\item \textbf{Filter Robustness:} The probability of successfully answering a Filter task decays linearly with the compression ratio $r$:
\begin{equation}
P(\text{Success}_{\mathcal{F}}) \geq 1 - \mathcal{O}((1-\rho) r)
\end{equation}
\item \textbf{Extract Collapse:} The probability of successfully answering an Extract task drops super-exponentially as a step-function dependent on the token rank. If the answer token's expected rank is $R$, then:
\begin{equation}
P(\text{Success}_{\mathcal{E}}) \approx \Phi\left( \frac{k - R}{\sigma} \right)
\end{equation}
where $\Phi$ is the standard normal CDF. As $k$ drops below $R$ (i.e., high compression $r$), the accuracy strictly collapses to $0$.
\end{enumerate}

\textit{Proof Sketch:}
For Filter Tasks, the presence of multiple high-relevance tokens means survival probability scales as $1 - r^m$, which is highly robust. For Extract Tasks, the unique target token has an expected rank $R$. When the retention threshold $k = n(1-r)$ decreases below $R$, the token is deterministically pruned, yielding $0.0$ accuracy.

\subsection*{A.3 The EA-KVzip Equivalence Principle}

\textbf{Motivation:} Expected Attention (EA) computes a statistical expectation in a single pass, while KVzip utilizes an expensive reconstruction mechanism. Yet, both achieve nearly identical F1 scores.

\textbf{Theorem A.3 (First-Order Taylor Equivalence):}
Expected Attention is the exact second-order Taylor approximation of the expected KVzip score, and the approximation error vanishes at a rate of $\mathcal{O}(d^{-3/2})$.

\textit{Proof:}
Let $f(q) = \exp(q \cdot k_i / \sqrt{d})$ be the unnormalized attention weight. We expand $f(q)$ around the mean future query $\mu_q$:
\begin{equation}
f(q) = f(\mu_q) + \nabla f(\mu_q)^T (q - \mu_q) + \frac{1}{2} (q - \mu_q)^T H(f)(\mu_q) (q - \mu_q) + \mathcal{O}(\|q-\mu_q\|^3)
\end{equation}
Taking the expectation over $q \sim \mathcal{N}(\mu_q, \Sigma_q)$:
\begin{equation}
\mathbb{E}[f(q)] = f(\mu_q) \left[ 1 + \frac{1}{2d} \text{Tr}(k_i k_i^T \Sigma_q) \right] = \exp\left( \frac{\mu_q \cdot k_i}{\sqrt{d}} \right) \left[ 1 + \frac{k_i^T \Sigma_q k_i}{2d} \right]
\end{equation}
For small $x$, $1 + x \approx e^x$. Therefore:
\begin{equation}
\mathbb{E}[f(q)] \approx \exp\left( \frac{\mu_q \cdot k_i}{\sqrt{d}} + \frac{k_i^T \Sigma_q k_i}{2d} \right)
\end{equation}
This matches the exact analytical formula used by EA. The Taylor remainder term scales as $\mathcal{O}(d^{-3/2})$, which becomes mathematically negligible for large hidden dimensions $d$.

\end{document}
